\par\Needspace{7\baselineskip}\bigskip
\begin{center}
\pdfbookmark[2]{Gauss an Encke, Göttingen, 25. März 1818}{pallas-brief-18180325}
\textsc{Gauss} an \textsc{Encke}, Göttingen, 25. März 1818.
\end{center}
\noindent\hspace*{\parindent}. . . . .\quad Ihre so oft erprobte Bereitwilligkeit, mir durch Ihre grosse Rechnungsfertigkeit Gefälligkeit zu erweisen, wage ich abermals in Anspruch zu nehmen, hoffend, dass der Gegenstand auch für Sie selbst nicht ohne Interesse sein werde. Meine Tafel für die Störungen der Pallas ist nunmehro ganz vollendet, ihre Bequemlichkeit lässt mir nichts zu wünschen übrig. Ich habe die sämmtlichen bisher beobachteten 12 Oppositionen von 1803 . . . . 1817 von neuem mit Hülfe der Tafel mit den Elementen verglichen, diese nachgefeilt und eine herrliche Übereinstimmung erreicht. Hiebei sind nun die so zahlreichen Beobachtungen von 1802, die aber erst nach der Opposition anfingen, gar nicht genutzt, und ich wäre sehr begierig zu wissen, welche Übereinstimmung meine Elemente in jenem Jahrgange geben; vielleicht wäre es doch thunlich, die Opposition daraus zu erhalten, da die Vergleichung mit Elementen geschehen kann, die der Wahrheit so nahe kommen müssen, dass \SourceLeaf{437}auch etwas entferntere Beobachtungen gebraucht werden dürften. Möchten Sie also nicht die Güte haben:
\begin{quote}
»die scharfe Vergleichung der sämmtlichen Meridianbeobachtungen\\
»der Pallas vom April und Mai 1802 mit den osculirenden Elementen\\
»zu übernehmen?«
\end{quote}
\par Die Elemente selbst sind zufolge meiner Tafeln folgende, wobei Knoten und Perihel sich auf das mittlere Äquinoctium von der Epoche beziehen und als siderisch ruhend zu betrachten sind:
\begin{center}
\begingroup\small\setlength{\tabcolsep}{4pt}
\begin{tabular}{@{}lr@{}}
\multicolumn{2}{c}{1802 April 19,49214 Göttinger Zeit}\\[3pt]
Mittlere Länge . . . . . . . . . . . . . .&\(166^0\;44'\;30''{,}43\)\\
Sonnennähe . . . . . . . . . . . . . . . . .&\(121\;12\;17{,}28\)\\
Aufsteigender Knoten . . . . . . . .&\(172\;27\;35{,}19\)\\
Excentricitätswinkel . . . . . . . . . .&\(14\;13\;27{,}83\)\\
Neigung der Bahn . . . . . . . . . . .&\(34\;37\;47{,}88\)\\
Tägl. mittl. tropische Bew. . . . .&\(770''{,}47396\)\\
Logarithm der halben Axe . . . .&\(0{,}442\,2173.\)
\end{tabular}
\endgroup
\end{center}
\noindent Die Epoche der Länge mögen Sie erst so corrigiren, dass\quad . . . . .
\par
